% Chapter 2: Literature Review
\chapter{Literature Review}
\label{ch:literature-review}

This chapter provides a comprehensive review of the existing literature. We build upon the foundational work mentioned in Chapter~\ref{ch:introduction}, particularly the contributions by Johnson~\cite{textbook2022} and recent advances~\cite{sample2023}.

\lipsum[9-10]

\section{Related Work}
\label{sec:related-work}

Several researchers have contributed to this area. The conference proceedings~\cite{conference2023} provide excellent examples of current approaches, while doctoral research~\cite{thesis2022} offers deeper theoretical insights.

\lipsum[11-12]

\subsection{Previous Studies}
\label{subsec:previous-studies}

\lipsum[13]

\section{Theoretical Framework}
\label{sec:theoretical-framework}

Building on the results from Chapter~\ref{ch:introduction}, we can now establish our theoretical framework.

\begin{block}[note]
The literature review process involved examining over 150 papers from the last decade to ensure comprehensive coverage of recent developments.
\end{block}

\begin{block}[tldr]
Key findings: Current methodologies show promise but lack comprehensive validation in real-world scenarios.
\end{block}

\begin{theorem}[Extension Theorem]
\label{thm:extension}
The results from Theorem~\ref{thm:main-result} can be extended to more general cases.
\end{theorem}

\begin{proof}
This follows by applying Lemma~\ref{lem:key-lemma} to the extended setting.
\end{proof}

\begin{example}[Concrete Application]
\label{ex:concrete}
Consider the case where our methodology is applied to real-world data. The theoretical guarantees still hold.
\end{example}

\lipsum[14-15]

\subsection{Key Concepts}
\label{subsec:key-concepts}

\lipsum[16]

\section{Gap Analysis}
\label{sec:gap-analysis}

\lipsum[17-18]